\subsection{Creusot}

Creusot is a \textit{deductive} verification tool for Rust whose development has started in 2020. The verifying process is semi-automatic: Creusot automatically generates theories from a given Rust code extended with Creusot annotations; these theories have to be imported into Why3 and evaluated there \cite{riesco_creusot_2022}. Creusot extends the idea of \textit{prophecies}, which have been first introduced to Rust by Matsushita et al. \cite{matsushita_rusthorn_2021, riesco_creusot_2022}. Creusot uses a language called \textit{Pearlite} for designing its specifications, which represents a pure subset of the Rust programming language extended with keywords and contains additional mathematical data types like Seq<T> and Int \cite{xavier_denis_xldeniscreusot_2023,riesco_creusot_2022}.

\subsubsection{Usability}

Like Prusti (refer to Chapter \ref{cap:prusti_usability}), Creusot integrates itself into the Rust compiler and uses the generated MIR for its theory generation \cite{riesco_creusot_2022}. The automatically generated theories produces from Creusot have to be manually imported into Why3, although there exists a convenient wrapper for that purpose \cite{riesco_creusot_2022, xavier_denis_xldeniscreusot_2023}.
From a developer standpoint, that results in additional effort and time required to import the generated theories into Why3 and evaluate them there. It also requires domain knowledge for formal verification and knowledge for using the Why3 ecosystem. 
From a formal verification expert standpoint, it opens a broad field of customization for specific needs e.g. setting a different timeout for the SMT solver for one particular theory.
It is planed to improve this part of the process though, in form of a Visual Studio Code plugin which acts as a frontend for Why3 \cite{xavier_denis_xldeniswhycode_2023}.

The annotations written in Pearlite are quite similar to those in Prusti, using the same Rust specific attributes style \cite{rust_foundation_attributes_nodate}, as can be seen in \autoref{lst:creusot_pre_postcondition}. Note that numbers in the specifications have to be equipped with the desired datatype for arithmetic operations, as Creusot automatically infers them as a Pearlite specific \textit{Int} datatype. Also note the practical appliance of the prophecy logic, which is annotated by the \textsuperscript{$\wedge$} operator in the postcondition and implies the final value of \textit{val}. 

It is also possible to call functions inside of the specifications like for example predicates, marked with the \lstinline{#[predicate]} annotation. Predicates are typically functions expecting Pearlite data types, preventing them being from called outside of verification purposes \cite{riesco_creusot_2022}. However, it is also possible to design helper predicate functions with original Rust data types. 

\begin{lstlisting}[language=Rust, caption=Example Rust code of a function with a pre- and postcondition for Creusot, label={lst:creusot_pre_postcondition}]
#[requires(*val <= u32::MAX - 1u32)]
#[ensures(^val == *val + 1u32)]
fn incr(val: &mut u32) {
    *val += 1;
}
\end{lstlisting}

\subsubsection{Technology}
%WhyML, MLCFG, Why3 -> Interchangeable SMT solver
%Macht keine Core proofs wie prusti also keine doppelte rust type checkung, dadurch schneller laut Paper
\subsubsection{Feedback}
% Why3 und Compile errors
\subsubsection{Language support and verification conditions}
% Safe Rust only, gibt auch #trusted für externe calls(?)
\subsubsection{Development}
\subsubsection{Activity}
\subsubsection{Scientific activity}
\subsubsection{Documentation}
% Testsuite mit praktischen Beispielen vorhanden, 